\section*{Chapter \ref{ch:b3:remote-sensing-gis} --- Remote Sensing and GIS}

\begin{solution}{exo:b3:remote-sensing-gis:1}
Field: $(0.42 - 0.06)/(0.42 + 0.06) = 0.75$, dense green vegetation, a crop
in full leaf. Ploughed field: $(0.28 - 0.21)/(0.28 + 0.21) = 0.14$, nearly
bare ground, a soil that brightens slightly towards the infrared.
\end{solution}

\begin{solution}{exo:b3:remote-sensing-gis:2}
(a)~$185000/30 \approx 6170$ pixels; $185000/15 \approx 12300$ in the
panchromatic band. (b)~$2^{12} = 4096$ levels; $2^8 = 256$. (c)~16 days
alone; with Landsat~9 in the same orbit 8 days apart, 8 days.
\end{solution}

\begin{solution}{exo:b3:remote-sensing-gis:3}
(a)~One fringe is $\lambda/2 = \qty{27.7}{mm}$; six fringes are
\qty{166}{mm} of line-of-sight change, here away from the satellite, since
the basin subsides. (b)~$56.6/2 = \qty{28.3}{mm}$.
\end{solution}

\begin{solution}{exo:b3:remote-sensing-gis:4}
(a)~$t = 2d/c = 2\times\num{5e5}/\num{3.00e8} = \qty{3.34}{ms}$.
(b)~$c\,\delta t/2 = \num{3.00e8}\times\num{8e-10}/2 = \qty{0.12}{m}$; a
surface sampled by many photons is located much better, the precision
falling as one over the square root of their number.
\end{solution}

\begin{solution}{exo:b3:remote-sensing-gis:5}
(a)~In SI units $L = \qty{9.0e6}{W.m^{-2}.sr^{-1}.m^{-1}}$. With $c_1 =
2hc^2 = \qty{1.191e-16}{W.m^2.sr^{-1}}$, $c_2 = \qty{0.014388}{m.K}$ and
$\lambda^5 = (\num{1.09e-5})^5 = \num{1.539e-25}$, $c_1/(\lambda^5L) = 86.0$,
and $T_b = c_2/(\lambda\ln 87.0) = \qty{295.6}{K}$. (b)~The kinetic temperature
solves $0.98\,B_\lambda(T) = L$, that is $B_\lambda(T) = L/0.98$: the same
formula gives $T = \qty{296.9}{K}$, \qty{1.3}{K} more.
\end{solution}

\begin{solution}{exo:b3:remote-sensing-gis:6}
(a)~$R = 693/\cos\ang{35} = \qty{846}{km}$. (b)~Real antenna: $R\lambda/L =
\num{8.46e5}\times0.0555/12.3 = \qty{3.8}{km}$; synthetic aperture: $L/2 =
\qty{6.2}{m}$, some 600 times finer. (c)~The synthetic aperture is the
length of track over which a point stays in the beam, $R\lambda/L =
\qty{3.8}{km}$.
\end{solution}

\begin{solution}{exo:b3:remote-sensing-gis:7}
$\partial z/\partial x = (418 - 405)/60 = 0.217$ (east minus west),
$\partial z/\partial y = (418 - 401)/60 = 0.283$ (north minus south);
$|\nabla z| = 0.357$, a slope of $\arctan 0.357 = \ang{19.6}$. The ground
rises to the north-east, so it faces the south-west: $A =
\operatorname{atan2}(-0.217, -0.283) = \ang{217}$. Hillshade: $\cos\ang{45}
\cos\ang{19.6} + \sin\ang{45}\sin\ang{19.6}\cos(\ang{315} - \ang{217}) =
0.666 - 0.033 = 0.63$, a little darker than flat ground (0.71): the slope
turns slightly away from a light in the north-west.
\end{solution}

\begin{solution}{exo:b3:remote-sensing-gis:8}
Row totals 52, 48, 38, 22; column totals 54, 46, 38, 22; diagonal 133 of
160: overall accuracy \qty{83.1}{\percent}. User's accuracies: water
$48/52 = 0.92$, forest $37/48 = 0.77$, crops $29/38 = 0.76$, town
$19/22 = 0.86$. Producer's: water $48/54 = 0.89$, forest $37/46 = 0.80$,
crops $29/38 = 0.76$, town $19/22 = 0.86$. $p_e = (52\cdot54 + 48\cdot46 +
38\cdot38 + 22\cdot22)/160^2 = 0.271$, $\kappa = (0.831 - 0.271)/(1 -
0.271) = 0.77$. Crops are most often missed: 9 of the 38 true crop pixels
went elsewhere, mostly to forest.
\end{solution}

\begin{solution}{exo:b3:remote-sensing-gis:9}
(a)~$L = 0.04\,B_\lambda(1443) + 0.96\,B_\lambda(290)$ at
\qty{10.9}{\micro\metre}, \qty{28.6}{W.m^{-2}.sr^{-1}.\micro\metre^{-1}},
whose \omterm{def:b3:remote-sensing-gis:brightness}{brightness temperature} is \qty{396}{K}, against a mean temperature of
$0.04\times1443 + 0.96\times290 = \qty{336}{K}$. (b)~At
\qty{2.2}{\micro\metre}, $B_\lambda(1443)/B_\lambda(290) = \num{6.8e7}$:
the lava dominates the pixel completely, and its \omterm{def:b3:remote-sensing-gis:brightness}{brightness temperature} is
\qty{845}{K}. The shortwave infrared band detects small hot spots far
better; the thermal band, sensitive to the whole pixel, is the one that
measures the ground around them.
\end{solution}

\begin{solution}{exo:b3:remote-sensing-gis:10}
With the satellite to the east, $d_1 = u_E\sin\theta + u_z\cos\theta$; to
the west, $d_2 = -u_E\sin\theta + u_z\cos\theta$. Adding and subtracting,
$u_z = (d_1 + d_2)/(2\cos\theta) = 102/1.638 = \qty{62.3}{mm}$ and $u_E =
(d_1 - d_2)/(2\sin\theta) = 26/1.147 = \qty{22.7}{mm}$. One pass gives one
equation for two unknowns: any mixture of uplift and eastward motion with
the same projection on the \omterm{def:b3:remote-sensing-gis:insar}{line of sight} gives the same fringes.
\end{solution}

\begin{solution}{exo:b3:remote-sensing-gis:11}
(a)~By \cref{thm:b3:remote-sensing-gis:fringe}, $\Delta\phi =
4\pi\Delta r/\lambda < \pi$ requires $\Delta r < \lambda/4 =
\qty{13.9}{mm}$. (b)~$u_z = k z_0(r^2 + z_0^2)^{-3/2}$ with $k = (1 -
\nu)\Delta V/\pi$; $\mathrm du_z/\mathrm dr = -3kz_0r(r^2 + z_0^2)^{-5/2}$,
whose derivative vanishes where $(r^2 + z_0^2) = 5r^2$, at $r = z_0/2$. There
\[
\left|\frac{\mathrm du_z}{\mathrm dr}\right| = \frac{3k}{2z_0^3}\Big(\frac54\Big)^{-5/2}
= \frac{0.859\,k}{z_0^3} = \frac{0.859\times\num{2.39e6}}{\num{6.4e10}} = \num{3.2e-5},
\]
\qty{32}{mm} per kilometre. (c)~$13.9/32 = \qty{0.43}{km}$: pixels smaller
than about \qty{430}{m}, easily met by Sentinel-1's.
\end{solution}

\begin{solution}{exo:b3:remote-sensing-gis:12}
(a)~For $N, R \ge 0$ not both zero, $|N - R| \le N + R$, so $|\mathrm{NDVI}|
\le 1$; replacing $N$ and $R$ by $aN$ and $aR$ multiplies the numerator and
the denominator by $a$. (b)~Without: $(0.502 - 0.051)/(0.502 + 0.051) =
0.816$; with: $(0.512 - 0.081)/(0.512 + 0.081) = 0.727$. (c)~The path
radiance lowers the index of green vegetation, by about a tenth here, and
by an amount that changes with the haze: a map of uncorrected images would
show the weather as much as the vegetation, and comparisons between dates
would be meaningless.
\end{solution}

\begin{solution}{pb:b3:remote-sensing-gis:1}
\textbf{1.} Forest $(0.38 - 0.04)/0.42 = 0.81$; flow of 1890 $0.23/0.37 =
0.62$; of 1960 $0.10/0.32 = 0.31$; of 2005 $0.02/0.26 = 0.08$.

\textbf{2.} 1890, then 1960, then 2005: the oldest flow is the greenest,
in the order of the ages; it has not yet reached the forest.

\textbf{3.} Bare rock reflects slightly more in the near infrared than in
the red, like the brown sand of the chapter's figure; lichens and the first
plants add a little: 0.08 is the value of nearly bare ground.

\textbf{4.} A slope turned away from the light dims both bands in the same
proportion; the ratio cancels it, while the near-infrared reflectance alone
would mistake a shaded forest for sparse vegetation.

\textbf{5.} By \cref{prop:b3:remote-sensing-gis:planck-inversion} with
$\lambda = \qty{10.9}{\micro\metre}$: the plateau \qty{291.1}{K}, the
fumarole pixel \qty{302.3}{K}.

\textbf{6.} $0.97\,B_\lambda(T) = \num{8.40e6}$ gives $T = \qty{293.0}{K}$.

\textbf{7.} $B_\lambda(\qty{370}{K}) =
\qty{22.48}{W.m^{-2}.sr^{-1}.\micro\metre^{-1}}$; $f = (9.95 - 8.40)/(22.48
- 8.40) = 0.110$, about \qty{1100}{m^2} of the \qty{e4}{m^2} pixel.

\textbf{8.} The radiance is a convex function of the temperature
(\cref{prop:b3:remote-sensing-gis:mixed}): the hot part adds more radiance
than its share of the temperature, so the pixel looks hotter than its mean,
$0.110\times370 + 0.890\times291.1 = \qty{299.8}{K}$ against
\qty{302.3}{K}.

\textbf{9.} The phase of each echo depends on the exact path; a large
change of viewpoint changes the sum of the scatterers' echoes and the
\omterm{def:b3:remote-sensing-gis:insar}{interferogram} loses its coherence, and it adds a phase from the relief.
$48/12 = 4$ cycles of one satellite.

\textbf{10.} $5.4\times27.73 = \qty{149.7}{mm}$ towards the satellite.

\textbf{11.} $u_z = 149.7/\cos\ang{35} = \qty{182.8}{mm}$.

\textbf{12.} The ground moves outwards as well as up; on the east side the
outward motion is towards the satellite and adds to the uplift, on the west
side it subtracts (\cref{prop:b3:remote-sensing-gis:mogi-los}). The change
falls faster to the west, where the fringes are crowded; the pattern is
stretched to the east.

\textbf{13.} On that line the horizontal motion is north--south,
perpendicular to the look, so the radar sees $u_z\cos\theta$ alone, of the
shape of the uplift. The uplift halves at $r = 0.766\,z_0$, so $z_0 =
3.4/0.766 = \qty{4.44}{km}$.

\textbf{14.} Half of 5.4: 2.7 fringes.

\textbf{15.} A deformation is the same on the two \omterm{def:b3:remote-sensing-gis:insar}{interferograms} after
scaling by the time intervals, or keeps growing from pair to pair; an
atmospheric pattern changes from date to date and often follows the
relief. A second pair that shows the same bull's eye confirms the
deformation.

\textbf{16.} $\Delta V = \pi u_z(0)z_0^2/(1 - \nu) = \pi\times0.1828\times
4436^2/0.75 = \qty{1.51e7}{m^3}$.

\textbf{17.} $\num{1.51e7}/48 = \qty{3.1e5}{m^3}$ a day, \qty{1.1e8}{m^3}
a year if the rate held.

\textbf{18.} $u_r \propto r/(r^2 + z_0^2)^{3/2}$ is largest where its
derivative vanishes, $r^2 + z_0^2 = 3r^2$, $r = z_0/\sqrt2 = \qty{3.14}{km}$;
there $u_r = \qty{70}{mm}$, \qty{38}{\percent} of the central uplift.

\textbf{19.} Radius \qty{1.48}{km}, volume $\frac43\pi(1.48)^3 =
\qty{13.6}{km^3}$; the change is $\num{1.51e7}/\num{1.36e10} =
\qty{0.11}{\percent}$ of it.

\textbf{20.} No: a tenth of a fringe is \qty{2}{\percent} of the count, and
the depth, read from a half-width of \qty{3.4}{km} perhaps to
\qty{0.1}{km}, enters squared, \qty{6}{\percent}; the \qty{3}{\percent}
correction is within that uncertainty.

\textbf{21.} The warming of the fumarole field (Part~II) over time, an
increase in small earthquakes beneath the bulge, a change in the gas
released, and uplift measured by satellite positioning or levelling at a
few points of the flank.

\textbf{22.} The reservoir under Mount Halvren grew by
\textbf{15~million~\unit[text-series-to-math=true]{m^3}} ($\num{1.51e7}\,\unit{m^3}$).
\end{solution}
